\documentclass[11pt,a4paper]{article}
\usepackage[utf8]{inputenc}
\usepackage[T1]{fontenc}
\IfFileExists{french.ldf}{\usepackage[french]{babel}}{\usepackage[english]{babel}}
\usepackage{amsmath,amssymb,amsthm}
\usepackage{geometry}\geometry{margin=2.3cm}
\usepackage{listings}
\usepackage{xcolor}
\usepackage{enumitem}
\usepackage{fancyhdr}
\usepackage{lmodern}
\usepackage{titlesec}
\usepackage{hyperref}
\hypersetup{colorlinks=true, urlcolor=[RGB]{2,101,115}, linkcolor=black}
\definecolor{codebg}{RGB}{247,249,251}
\definecolor{kw}{RGB}{175,45,45}
\definecolor{cmt}{RGB}{110,120,130}
\definecolor{str}{RGB}{20,110,60}
\lstset{language=Python, backgroundcolor=\color{codebg}, basicstyle=\ttfamily\small,
  keywordstyle=\color{kw}\bfseries, commentstyle=\color{cmt}\itshape, stringstyle=\color{str},
  numbers=left, numberstyle=\tiny\color{cmt}, numbersep=8pt, frame=single, rulecolor=\color{gray!40},
  showstringspaces=false, breaklines=true, columns=fullflexible, extendedchars=true, inputencoding=utf8,
  literate={é}{{\'e}}1 {è}{{\`e}}1 {à}{{\`a}}1 {ê}{{\^e}}1 {î}{{\^i}}1 {ç}{{\c c}}1 {ô}{{\^o}}1 {û}{{\^u}}1 {ù}{{\`u}}1 {É}{{\'E}}1}
\newcounter{qcnt}
\newcommand{\Q}{\medskip\par\stepcounter{qcnt}\noindent\textbf{Question~\theqcnt.}~}
\newcommand{\code}[1]{\lstinline!#1!}
\pagestyle{fancy}\fancyhf{}
\lhead{\small Préparation au CNC 2026 --- Informatique MP}
\rhead{\small Arbres de Merkle et preuves d'inclusion}
\cfoot{\thepage}
\renewcommand{\headrulewidth}{0.4pt}
\titleformat*{\section}{\large\bfseries\scshape}

\begin{document}

\begin{center}
{\large\bfseries Préparation au Concours National Commun 2026 --- Informatique MP}\\[8pt]
\rule{0.9\linewidth}{0.5pt}\\[6pt]
{\Large\bfseries Épreuve d'Informatique}\\[4pt]
{\bfseries Filière MP --- Durée : 4 heures}\\[6pt]
{\itshape Arbres de Merkle et preuves d'inclusion : hachage, construction récursive, certificats logarithmiques}\\[6pt]
\rule{0.9\linewidth}{0.5pt}\\[8pt]
{\small Auteur~: \href{https://orcid.org/0000-0002-6108-7176}{\textbf{Pr Agrégé El Hadiq Zouhair}}}\\[3pt]
{\small \href{https://cpgeacademy.org}{\textbf{cpgeacademy.org}} \quad---\quad Tous droits réservés \textcopyright\ cpgeacademy.org}
\end{center}
\medskip
\begin{center}
\fcolorbox{gray}{gray!5}{\begin{minipage}{0.93\linewidth}\small
\textbf{Avis important.} Ce document est une \textbf{initiative personnelle} du professeur,
à but strictement pédagogique. Il ne s'agit \textbf{pas} d'un sujet officiel du Concours
National Commun~; il s'en inspire uniquement par la forme et l'esprit.
\smallskip
\textbf{Pour aller plus loin.} Les élèves peuvent traiter une \textbf{nouvelle version} de
ce problème dotée d'un \textbf{autocorrecteur des réponses}, et suivre un \textbf{cours
complet}, sur \href{https://cpgeacademy.org}{\textbf{cpgeacademy.org}}.
\end{minipage}}
\end{center}

\bigskip
\section*{Présentation et notations}

Un \emph{arbre de Merkle} permet de résumer une grande collection de données par un unique
condensé, la \emph{racine}, et de prouver qu'une donnée en fait partie à l'aide d'un
\emph{certificat} de taille logarithmique. C'est un outil central des blockchains (preuves
d'inclusion, listes d'autorisation, stockage vérifiable). Dans tout le problème, les
\emph{feuilles} sont des entiers, et l'on combine deux n\oe uds par la fonction
$$\mathtt{parent}(a,b) = (31\,a + b) \bmod p, \qquad p = 1000003 \ (\text{premier}).$$

\noindent\textbf{Conventions.} Le langage est Python 3. Pour une liste \code{L}, \code{len(L)}
est sa longueur, \code{L[i]} son terme d'indice $i$ ($\ge 0$), \code{L[-1]} son dernier terme.
$\lceil x \rceil$ est la partie entière supérieure, $\log_2$ le logarithme binaire. Les questions
sont numérotées \textbf{en continu}.

\bigskip
\hrule
\bigskip
{\Large\bfseries É\,N\,O\,N\,C\,É}
\medskip

\section*{Partie I --- La fonction de combinaison}

\Q Écrire \code{parent(a, b)} conforme à la définition ci-dessus.

\Q Calculer à la main $\mathtt{parent}(1,2)$, $\mathtt{parent}(3,4)$ et
$\mathtt{parent}(33,97)$.

\Q On dit que \code{parent} est \emph{déterministe}. Expliquer pourquoi ce caractère déterministe
est indispensable pour que deux n\oe uds différents du réseau, partant des mêmes feuilles,
calculent la même racine.

\Q On admet, pour la suite, que \code{parent} est \emph{résistante aux collisions}~: on ne sait
pas exhiber deux couples distincts $(a,b)\neq(a',b')$ tels que $\mathtt{parent}(a,b)=
\mathtt{parent}(a',b')$. Rappeler pourquoi une fonction $h : \mathbb{N}^2 \to \{0,\dots,p-1\}$ ne
peut pas être \emph{injective}, et donc pourquoi la résistance aux collisions est une hypothèse
\emph{calculatoire} (« difficile à violer ») et non une impossibilité mathématique.

\section*{Partie II --- Construction de la racine}

La racine se calcule par niveaux~: tant qu'il reste plus d'un n\oe ud, si le nombre de n\oe uds
est impair on \textbf{duplique le dernier}, puis on remplace la liste par celle des
\code{parent} de couples consécutifs.

\Q Écrire \code{racine_merkle(feuilles)} (liste non vide) renvoyant la racine.

\Q Dérouler le calcul de \code{racine_merkle([1, 2, 3, 4])} niveau par niveau et donner la valeur
de la racine.

\Q Faire de même pour \code{racine_merkle([1, 2, 3])} en explicitant l'étape de duplication.

\Q Soit $n$ le nombre de feuilles. En notant $t_0=n$ et $t_{k+1}=\lceil t_k/2\rceil$ la taille du
niveau $k+1$, montrer que le nombre de niveaux est $O(\log n)$ et que le nombre total d'appels à
\code{parent} est un $O(n)$.

\section*{Partie III --- Preuves d'inclusion}

\Q Écrire \code{preuve_inclusion(feuilles, i)} renvoyant la liste des couples
$(\text{frère}, \text{côté})$, avec côté valant \code{'L'} ou \code{'R'}, rencontrés en remontant
de la feuille d'indice $i$ jusqu'à la racine.

\Q Écrire \code{verifier(feuille, preuve, racine)} qui recalcule la racine depuis la feuille et
la preuve et renvoie un booléen. Vérifier sur l'exemple $[1,2,3,4]$, $i=0$.

\Q Montrer que la taille d'une preuve d'inclusion est le nombre de niveaux, soit $O(\log n)$, et
commenter l'intérêt (comparer au coût de transmettre toutes les feuilles).

\Q \textbf{Soundness.} Démontrer par récurrence sur la hauteur $h$ de l'arbre que, sous
l'hypothèse de résistance aux collisions, si \code{verifier(feuille, preuve, racine)} renvoie
\code{True} alors \code{feuille} figure effectivement parmi les feuilles ayant produit
\code{racine}.

\section*{Partie IV --- Hauteur et mise à jour}

\Q Montrer que la hauteur de l'arbre (nombre de niveaux au-dessus des feuilles) vaut
$H(n) = \lceil \log_2 n \rceil$ pour $n \ge 1$ (avec $H(1)=0$). Donner $H(n)$ pour
$n \in \{1,2,3,4,5,8,9\}$.

\Q Lorsqu'une seule feuille change de valeur, quels n\oe uds de l'arbre doivent être recalculés~?
En déduire qu'une mise à jour de la racine coûte $O(\log n)$ appels à \code{parent} (et non
$O(n)$). Écrire \code{maj_racine(feuilles, i, v)} qui renvoie la nouvelle racine lorsque la feuille
d'indice $i$ prend la valeur \code{v}, en \textbf{réutilisant} \code{racine_merkle}. \emph{On ne
demande pas ici l'implémentation optimale en $O(\log n)$, mais on justifiera la borne.}

\newpage
\setcounter{qcnt}{0}
\bigskip
\hrule
\bigskip
{\Large\bfseries C\,O\,R\,R\,I\,G\,É}
\medskip

\section*{Partie I --- La fonction de combinaison}

\Q
\begin{lstlisting}
def parent(a, b):
    return (31 * a + b) % 1000003
\end{lstlisting}

\Q $\mathtt{parent}(1,2) = 31\cdot1+2 = 33$~; $\mathtt{parent}(3,4)=31\cdot3+4=97$~;
$\mathtt{parent}(33,97)=31\cdot33+97=1023+97=1120$.

\Q Le calcul de la racine doit donner \emph{le même} résultat sur tous les n\oe uds du réseau
pour que le consensus soit possible. Si \code{parent} pouvait renvoyer des valeurs différentes
pour les mêmes entrées (aléa, dépendance à l'environnement), deux n\oe uds honnêtes partant des
mêmes feuilles obtiendraient des racines différentes et ne pourraient pas s'accorder. Le
déterminisme garantit la \emph{reproductibilité}.

\Q L'ensemble de départ $\mathbb{N}^2$ est infini (ou, en pratique, de cardinal bien plus grand
que $p$), tandis que l'ensemble d'arrivée $\{0,\dots,p-1\}$ est de cardinal fini $p$. Par le
principe des tiroirs, $h$ ne peut être injective~: des collisions \emph{existent}
nécessairement. La résistance aux collisions n'affirme donc pas leur absence, mais qu'il est
\emph{calculatoirement difficile} d'en exhiber une~: c'est une hypothèse de sécurité, pas un
théorème d'injectivité. $\square$

\section*{Partie II --- Construction de la racine}

\Q
\begin{lstlisting}
def racine_merkle(feuilles):
    niveau = list(feuilles)
    while len(niveau) > 1:
        if len(niveau) % 2 == 1:
            niveau.append(niveau[-1])       # duplique le dernier
        niveau = [parent(niveau[i], niveau[i+1])
                  for i in range(0, len(niveau), 2)]
    return niveau[0]
\end{lstlisting}

\Q Niveau $0=[1,2,3,4]$~; $\mathtt{parent}(1,2)=33$, $\mathtt{parent}(3,4)=97$, donc niveau
$1=[33,97]$~; $\mathtt{parent}(33,97)=1120$. Racine $=\mathbf{1120}$.

\Q Niveau $0=[1,2,3]$, impair $\Rightarrow$ on duplique~: $[1,2,3,3]$. $\mathtt{parent}(1,2)=33$,
$\mathtt{parent}(3,3)=31\cdot3+3=96$, donc niveau $1=[33,96]$~; $\mathtt{parent}(33,96)=1023+96=
1119$. Racine $=\mathbf{1119}$.

\Q On a $t_{k+1}=\lceil t_k/2\rceil \le t_k/2 + 1$. Par récurrence $t_k \le n/2^{k} + 2$, donc
$t_k \le 2$ dès que $n/2^{k}\le 1$, c'est-à-dire $k \ge \log_2 n$~: il y a $O(\log n)$ niveaux. Le
nombre d'appels à \code{parent} au niveau $k$ est $\lceil t_k/2\rceil$, et le total vaut
$\sum_{k\ge 1} t_k \le \sum_{k\ge1}(n/2^k + 2) \le n + O(\log n) = O(n)$. $\square$

\section*{Partie III --- Preuves d'inclusion}

\Q
\begin{lstlisting}
def preuve_inclusion(feuilles, i):
    niveau = list(feuilles)
    preuve = []
    idx = i
    while len(niveau) > 1:
        if len(niveau) % 2 == 1:
            niveau.append(niveau[-1])
        if idx % 2 == 0:
            preuve.append((niveau[idx+1], 'R'))   # frere a droite
        else:
            preuve.append((niveau[idx-1], 'L'))   # frere a gauche
        niveau = [parent(niveau[j], niveau[j+1])
                  for j in range(0, len(niveau), 2)]
        idx //= 2
    return preuve
\end{lstlisting}

\Q
\begin{lstlisting}
def verifier(feuille, preuve, racine):
    acc = feuille
    for frere, cote in preuve:
        if cote == 'R':
            acc = parent(acc, frere)
        else:
            acc = parent(frere, acc)
    return acc == racine
\end{lstlisting}
Exemple~: $\mathtt{preuve\_inclusion}([1,2,3,4],0)=[(2,\mathtt{'R'}),(97,\mathtt{'R'})]$. La
vérification calcule $\mathtt{parent}(1,2)=33$ puis $\mathtt{parent}(33,97)=1120$, égal à la
racine~: \code{True}.

\Q La preuve contient exactement un couple par niveau au-dessus de la feuille, soit $H(n)=
\lceil\log_2 n\rceil = O(\log n)$ couples. On transmet donc $O(\log n)$ valeurs au lieu des $n$
feuilles~: pour $n=10^{6}$, une vingtaine de valeurs suffisent. C'est ce qui rend les preuves de
Merkle utilisables sur une blockchain, où l'espace de stockage on-chain est coûteux.

\Q \textbf{Récurrence sur $h$.} \emph{Base} $h=0$~: l'arbre est réduit à une feuille, la racine
\emph{est} cette feuille, et la preuve est vide~; \code{verifier} renvoie \code{True} seulement si
$\mathtt{feuille}=\mathtt{racine}$, donc la feuille est bien celle ayant produit la racine.
\emph{Hérédité}~: supposons le résultat vrai pour tout arbre de hauteur $h$, et soit un arbre de
hauteur $h+1$. La vérification calcule d'abord $r' = \mathtt{parent}(\mathtt{feuille},f)$ (ou
$\mathtt{parent}(f,\mathtt{feuille})$), puis vérifie que $r'$, muni du reste de la preuve, mène à
la racine~: c'est une preuve d'inclusion valide de $r'$ dans l'arbre des n\oe uds de niveau $1$,
de hauteur $h$. Par hypothèse de récurrence, $r'$ est bien un n\oe ud de ce niveau, produit par
un couple réel de l'arbre. Comme \code{parent} résiste aux collisions, le seul antécédent
\emph{connu} de $r'$ est le couple effectivement présent~: c'est $(\mathtt{feuille},f)$. Donc
\code{feuille} figure parmi les feuilles ayant produit la racine. $\square$

\section*{Partie IV --- Hauteur et mise à jour}

\Q À chaque niveau, la taille passe de $t$ à $\lceil t/2\rceil$~; partant de $n$, il faut le plus
petit $H$ tel que $\lceil\cdots\lceil n/2\rceil\cdots\rceil = 1$ après $H$ divisions, soit
$2^{H-1} < n \le 2^{H}$ pour $n \ge 2$, d'où $H(n)=\lceil\log_2 n\rceil$ (et $H(1)=0$). On obtient
$$H(1)=0,\ H(2)=1,\ H(3)=2,\ H(4)=2,\ H(5)=3,\ H(8)=3,\ H(9)=4.$$

\Q Seuls les n\oe uds situés \emph{sur le chemin} de la feuille modifiée jusqu'à la racine
changent de valeur~: à chaque niveau, un seul n\oe ud dépend de la feuille modifiée (son
ancêtre). Il y a $H(n)=O(\log n)$ tels n\oe uds~; recalculer chacun coûte un appel à \code{parent},
d'où une mise à jour en $O(\log n)$ appels. Une implémentation simple, correcte mais non
optimale, recalcule toute la racine~:
\begin{lstlisting}
def maj_racine(feuilles, i, v):
    nouvelles = list(feuilles)
    nouvelles[i] = v
    return racine_merkle(nouvelles)
\end{lstlisting}
Cette version coûte $O(n)$~; l'argument ci-dessus montre qu'une version conservant les n\oe uds
inchangés atteindrait $O(\log n)$. $\square$

\medskip
\noindent\emph{Vérifications numériques} (exécutées en Python)~: $\mathtt{racine\_merkle}([1,2,3,4])
=1120$, $\mathtt{racine\_merkle}([1,2,3])=1119$, et les hauteurs annoncées correspondent à
$\lceil\log_2 n\rceil$.

\end{document}
